\section*{第 \ref{ch:b1:clastic-sediments} 章——碎屑沉积及其环境}

\begin{solution}{exo:b1:clastic-sediments:1}
$\varphi = -\log_2 d$：\qty{0.18}{mm} 得 $2.47$，为细砂（2 至 3）；
\qty{0.75}{mm} 得 $0.42$，为粗砂（0 至 1）；\qty{3}{mm} 得
$-1.58$，为砾石（$-1$ 至 $-2$）；\qty{0.012}{mm} 得 $6.38$，为细粉砂（6
至 7）；\qty{150}{mm} 得 $-7.23$，为卵石（$-6$ 至 $-8$）。
\end{solution}

\begin{solution}{exo:b1:clastic-sediments:2}
$M_z = (1.05 + 1.71 + 2.43)/3 = 1.73\,\varphi$，即 $2^{-1.73} =
\qty{0.30}{mm}$，为中砂。$\sigma_I = (2.43 - 1.05)/4 + (2.98 - 0.62)/6.6
= 0.345 + 0.358 = 0.70\,\varphi$：\omterm{def:b1:clastic-sediments:sorting}{分选}中等偏好（0.50 至 0.71），略
低于中等\omterm{def:b1:clastic-sediments:sorting}{分选}的界限。
\end{solution}

\begin{solution}{exo:b1:clastic-sediments:3}
$q_b = (1 - 0.4) \times 0.02 \times (\num{0.01}/60)/2 = \qty{1.0e-6}{m^2.s^{-1}}$
固体颗粒。一天内，$\num{1.0e-6} \times 86400 = \qty{0.0864}{m^3}$，每
米脊线，或 $0.0864 \times 2650 = \qty{229}{kg}$ 石英。
\end{solution}

\begin{solution}{exo:b1:clastic-sediments:4}
(a)~50 节为 $50 \times 1852/3600 = \qty{25.7}{m.s^{-1}}$；12 节，
\qty{6.17}{m.s^{-1}}。(b)~$\qty{100000}{m}/\qty{6.17}{m.s^{-1}} =
\qty{16200}{s}$，约 \qty{4.5}{h}。(c)~480 英里为 \qty{772}{km}，400 英里
为 \qty{644}{km}。
\end{solution}

\begin{solution}{exo:b1:clastic-sediments:5}
(a)~$(\varphi_{84} - \varphi_{16})/2 = 2 \times 0.9945\,\sigma/2 =
0.9945\,\sigma$，而 $(\varphi_{95} - \varphi_5)/3.3 = 2 \times
1.6449\,\sigma/3.3 = 0.9969\,\sigma$：两者对 $\sigma$ 的估计都在三分之一
个百分点以内。(b)~实际沉积物未必在 φ 标度下服从正态分布：少量
远更粗或远更细颗粒的混入会大幅改变 $\varphi_5$ 或
$\varphi_{95}$，而 $\varphi_{16}$ 和 $\varphi_{84}$ 变化很小，因此第二种估计值
增大，第一种则不变。两者在混合物中差异最大，例如
含若干砾石和泥的砂，或双峰沉积。
\end{solution}

\begin{solution}{exo:b1:clastic-sediments:6}
(a)~$\sin\theta = T/L = 1.5/15 = 0.1$，故 $\theta = \ang{5.74}$。(b)~$a =
c\tan\theta = 0.5 \times 0.1005 = \qty{0.050}{cm}$ 每分钟，或
\qty{3.0}{cm} 每小时。(c)~床面每小时抬升数厘米，说明其接收的沙子
正从悬浮状态迅速沉降：这是一股载沙很重、正在减速的水流，
例如消退中的洪水或\omterm{def:b1:clastic-sediments:turbidite}{浊流}的尾部。
\end{solution}

\begin{solution}{exo:b1:clastic-sediments:7}
(a)~$(350 + 20 + 5 + 340 + 15 + 30)/6 = 760/6 = 126.7$：东南方向，
尽管每个读数都在北偏 \ang{20} 以内。(b)~$S =
\sum\sin\alpha_i = 0.672$，$C = \sum\cos\alpha_i = 5.692$，故 $\bar\alpha =
\operatorname{atan2}(0.672, 5.692) = \ang{6.7}$，且 $\bar R =
\sqrt{0.672^2 + 5.692^2}/6 = 5.732/6 = 0.955$。(c)~算术平均把
350 与 005 视为相距很远，而在圆周上二者仅相距 \ang{15}：\omterm{def:b1:spherical-astronomy:horizontal}{方位角}须
作为单位矢量求平均。较大的 $\bar R$ 表明各读数
紧密集中在 \ang{7} 附近。
\end{solution}

\begin{solution}{exo:b1:clastic-sediments:8}
A 的黏土含量超过 \qty{5}{\percent}：未成熟，成分上
也未成熟，含\omterm{def:b1:rock-forming-minerals:groups}{长石}和岩屑；它快速沉积于靠近物源处，
形成于冲积扇或河漫滩，或由\omterm{def:b1:clastic-sediments:turbidite}{浊流}搬运而成。B 的黏土含量低于
\qty{5}{\percent}，\omterm{def:b1:clastic-sediments:sorting}{分选}良好（$\sigma_I < 0.5\,\varphi$），但颗粒
为次圆状（$\rho < 3$）：成熟，例如河流边滩，或陆架上
经水流淘洗的砂。C 不含黏土，\omterm{def:b1:clastic-sediments:sorting}{分选}极好，颗粒
为圆状（$\rho > 3$）：超成熟，石英含量达 \qty{98}{\percent}，
成分上也成熟，如海滩或沙漠\omterm{def:b1:clastic-sediments:bedform}{沙丘}区的沉积。
\end{solution}

\begin{solution}{exo:b1:clastic-sediments:9}
(a)~$1.5/\sqrt{9.81 \times 0.3} = 0.87$；$1.0/\sqrt{9.81 \times 2} = 0.23$。
(b)~$h = U^2/g = 1.5^2/9.81 = \qty{0.23}{m}$；$U = \sqrt{9.81 \times 2} =
\qty{4.43}{m.s^{-1}}$。(c)~在给定速度下，弗劳德数随水深减小而增大：
薄而快的水层，如海滩上波浪的回流，
很容易达到接近 1 的弗劳德数，而水深数米的河流
则需要每秒数米的流速，除洪水期外很少有河流能达到。
\end{solution}

\begin{solution}{exo:b1:clastic-sediments:10}
(a)~取中点 $-0.5, 0.5, \ldots, 4.5$：$m = (-1 + 4 + 45 + 100 + 52.5
+ 22.5)/100 = 2.23\,\varphi$；以百分比加权的平方偏差之和为
\num{107.71}，故 $s = \sqrt{1.0771} = 1.04\,\varphi$：
\omterm{def:b1:clastic-sediments:sorting}{分选}差。(b)~累积粗于该粒径的百分比依次为 0、2、10、40、80、95 和
100，对应 $-1, 0, \ldots, 5\,\varphi$。于是 $\varphi_5 = 0 + 3/8 = 0.375$，
$\varphi_{16} = 1 + 6/30 = 1.20$，$\varphi_{50} = 2 + 10/40 = 2.25$，
$\varphi_{84} = 3 + 4/15 = 3.27$，$\varphi_{95} = 4.0$。$M_z = (1.20 + 2.25 +
3.27)/3 = 2.24\,\varphi$，且 $\sigma_I = (3.27 - 1.20)/4 + (4.0 - 0.375)/6.6 =
0.517 + 0.549 = 1.07\,\varphi$。(c)~两种方法的结果一致：平均值相差约
\num{0.01}，\omterm{def:b1:clastic-sediments:sorting}{分选}相差约 \num{0.03}：
为细砂（$2.24\,\varphi$ 即 \qty{0.21}{mm}），\omterm{def:b1:clastic-sediments:sorting}{分选}差。
\end{solution}

\begin{solution}{exo:b1:clastic-sediments:11}
(a)~$U = \sqrt{0.02 \times 9.81 \times 50 \times 0.005/0.003} =
\qty{4.04}{m.s^{-1}}$，或 7.9 节。(b)~$(\Delta\rho/\rho)\,h = U^2
C_f/(gS) = 25.7^2 \times 0.003/(9.81 \times 0.005) = \qty{40}{m}$：例如
一股厚 \qty{400}{m}、密度超出 \qty{10}{\percent} 的\omterm{def:b1:clastic-sediments:turbidite}{浊流}，
由此可见 1929 年那股\omterm{def:b1:clastic-sediments:turbidite}{浊流}必定十分巨大且密度很高。(c)~流速
与 $\sqrt S$ 成正比：$4.04/\sqrt 5 = \qty{1.81}{m.s^{-1}}$。较慢的\omterm{def:b1:clastic-sediments:turbidite}{浊流}
其湍流较弱，会沉积较粗的颗粒；其密度变小，流速
进一步减慢，并以\omterm{def:b1:clastic-sediments:turbidite}{递变层}的形式将载荷铺布于平原之上。
\end{solution}

\begin{solution}{exo:b1:clastic-sediments:12}
(a)~$0.7 \times \qty{5e10}{kg}/\qty{2000}{kg.m^{-3}} = \qty{1.75e7}{m^3}$ 每
年。(b)~将其铺展于高 \qty{30}{m}、宽 \qty{100}{km} 的锋面上，
使锋面推进 $\num{1.75e7}/(30 \times \num{1e5}) = \qty{5.8}{m}$ 每
年。(c)~$a = v\tan\alpha = 5.83 \times \tan\ang{0.5} = \qty{0.051}{m}$，约
\qty{5}{cm} 每年，比推进速度小一百倍；
\qty{30}{m} 厚的层序需 $30/0.051 = 590$ 年才能形成。
\end{solution}

\begin{solution}{pb:b1:clastic-sediments:1}
\textbf{1.} 单元~1：$(5.0 + 6.6 + 8.4)/3 = 6.67\,\varphi$，\qty{0.0098}{mm}
(\qty{9.8}{\micro m})；单元~2：$3.47\,\varphi$，\qty{0.090}{mm}；单元~3：
$1.80\,\varphi$，\qty{0.29}{mm}；单元~4：$0.73\,\varphi$，\qty{0.60}{mm}。

\textbf{2.} 单元~1：$(8.4 - 5.0)/4 + (9.6 - 4.2)/6.6 = 0.850 + 0.818 =
1.67\,\varphi$；单元~2：$0.550 + 0.545 = 1.10$；单元~3：$0.275 + 0.273 = 0.55$；
单元~4：$0.725 + 0.712 = 1.44$。

\textbf{3.} 单元~1：细粉砂，\omterm{def:b1:clastic-sediments:sorting}{分选}差；单元~2：极细砂，
\omterm{def:b1:clastic-sediments:sorting}{分选}差；单元~3：中砂，\omterm{def:b1:clastic-sediments:sorting}{分选}中等偏好；单元~4：粗砂，
\omterm{def:b1:clastic-sediments:sorting}{分选}差。

\textbf{4.} 单元~3。河口处的水流和波浪反复搬动砂，并将细粒带走，
而前三角洲泥则由悬浮中沉降下来，
是抵达此处的各细粒级的混合物。

\textbf{5.} $8\,\varphi$ 位于 $1.4/1.8$ 处（自 $\varphi_{50}$ 至 $\varphi_{84}$ 计），
因此 $50 + (1.4/1.8) \times 34 = 76$，即样品的 76\% 较粗，
约 \qty{24}{\percent} 为黏土级。这样细的颗粒只能从近乎静止的水中沉降：
该处低于波浪与水流的
作用范围。

\textbf{6.} 河道既携带洪水时在河床上滚动的砾石，也携带砂；
沉积物将最深水流形成的砾石滞留层与较浅水流的砂相混合，
因而涵盖多个粒级。

\textbf{7.} \omterm{def:b1:clastic-sediments:sorting}{平均粒径}由 \qty{0.0098}{mm} 增至 \qty{0.090}{mm}，再至 \qty{0.29}{mm}：
该层序向上变粗。

\textbf{8.} 单元~1，纹层状泥：前三角洲；单元~2，粉砂与砂互层：
下三角洲前缘；单元~3，\omterm{def:b1:clastic-sediments:cross-bedding}{交错层理}砂：上三角洲前缘的
河口坝。

\textbf{9.} 每一带都覆盖于先前位于其前方的那一带之上：
三角洲向海方向持续推进，并覆盖于该地点之上。

\textbf{10.} 至少 \qty{21}{m}，即单元~1至3的厚度，因为该层序
充填了从前三角洲底部到近表河口坝之间的水体；
此后泥质已被压实，厚度还应更大
（见\cref{ch:b1:diagenesis-basins}）。

\textbf{11.} 三角洲平原上的一条分支河道，随三角洲向前推进而切入河口坝；
其卵石即为\omterm{def:b1:erosion-transport:load}{推移质}。

\textbf{12.} 不是：点坝以侵蚀面和滞留层开始，并向上变细，过渡为河漫滩泥。
三角洲前缘则由泥向上变粗。

\textbf{13.} $\qty{16}{m}/\num{4000}\text{ 年} = \qty{4.0}{mm}$ 每年。

\textbf{14.} $v = a/\tan\alpha = 0.0040/\tan\ang{0.5} = 0.0040/0.008727 =
\qty{0.458}{m}$ 每年。

\textbf{15.} $0.458 \times 4000 = \qty{1830}{m}$，约 \qty{1.8}{km}。

\textbf{16.} $\qty{21}{m}/\qty{4.0}{mm.\text{年}^{-1}} = 5250$ 年。

\textbf{17.} 三角洲前缘接收的砂大部分位于河口附近，前三角洲接收的则少得多，
三角洲平原除洪水期外所获甚少；活动的河口也会不断迁移，
因此某一点在一段时间内可能获得大量沉积，随后几乎不再接收。
两层定年层之间的单一速率，是对这些变化的平均，
而泥质因压实损失了部分厚度。

\textbf{18.} 总和为 1544，故算术平均为 \ang{128.7}，即东南向，
毫无意义，因为读数围绕北向散布，
而 340 附近与 010 附近的数值被当作相距甚远来取平均。

\textbf{19.} $S = \sum\sin\alpha_i = 1.698$ 且 $C = \sum\cos\alpha_i =
11.248$。

\textbf{20.} $\bar R = \sqrt{1.698^2 + 11.248^2}/12 = 11.376/12 = 0.948$：
各读数紧密集中于同一方向。

\textbf{21.} 河流流动，三角洲向大致北方推进
（见\cref{prop:b1:clastic-sediments:azimuths}）；陆地位于南侧。

\textbf{22.} 由于前积层是在已不再水平的层面上测量的：
其\omterm{def:b1:geological-maps:attitude}{倾向}为倾斜层的\omterm{def:b1:geological-maps:attitude}{倾向}加上倾斜作用。
将每个前积层绕层面\omterm{def:b1:geological-maps:attitude}{走向}旋转该层面的\omterm{def:b1:geological-maps:attitude}{倾角}，
层面即恢复水平，前积层也随之
恢复到原始\omterm{def:b1:geological-maps:attitude}{产状}。

\textbf{23.} 潮汐流往复流动，形成\omterm{def:b1:geological-maps:attitude}{倾向}相反的两组层系，如人字形\omterm{def:b1:clastic-sediments:cross-bedding}{交错层理}；
玫瑰图应有两个相隔约 \ang{180} 的瓣。
十二个读数构成一个瓣，
$\bar R$ 接近 1，表示河流水流朝一个方向流动。

\textbf{24.} $\bar\alpha = \operatorname{atan2}(1.698, 11.248) = \ang{8.6}$：
单元~3 的\omterm{def:b1:clastic-sediments:cross-bedding}{古水流}\omterm{def:b1:field-lab-methods-1:mean-orientation}{平均方位}角为 \ang{9}，北偏东少许。
\end{solution}
